\documentclass[11pt,a4paper]{article}

\usepackage[T1]{fontenc}
\usepackage[utf8]{inputenc}
\usepackage{lmodern}
\usepackage{amssymb,amsmath,amsthm,mathtools}
\usepackage[margin=2.5cm]{geometry}
\usepackage{parskip}
\usepackage[hidelinks]{hyperref}
\usepackage{microtype}

% Theorem environments (arXiv / journal style)
\theoremstyle{plain}
\newtheorem{theorem}{Theorem}[section]
\newtheorem{proposition}[theorem]{Proposition}
\newtheorem{lemma}[theorem]{Lemma}
\newtheorem{corollary}[theorem]{Corollary}

\theoremstyle{definition}
\newtheorem{definition}[theorem]{Definition}
\newtheorem{example}[theorem]{Example}
\newtheorem{remark}[theorem]{Remark}

% Notation
\newcommand{\Ext}{\mathrm{Ext}}
\newcommand{\Shape}{\mathrm{Shape}}
\newcommand{\Pos}{\mathrm{Pos}}
\newcommand{\Type}{\mathbf{Type}}
\newcommand{\id}{\mathrm{id}}
\newcommand{\rootpos}{\mathsf{o}}
\newcommand{\sub}{\mathbin{\downarrow}}
\newcommand{\shift}{\oplus}
\newcommand{\cat}{\mathcal{C}}
\newcommand{\DCont}{\mathbf{DCont}}
\newcommand{\Cof}{\mathbf{Cof}}
\newcommand{\Cat}{\mathbf{Cat}}
\newcommand{\Hom}{\mathrm{Hom}}

\title{Directed Containers as Categories\\
  and Their Morphisms as Cofunctors}
\author{Extracted from a machine-checked development}
\date{}

\begin{document}
\maketitle

\begin{abstract}
We formalise the object-level dictionary between directed containers and small
categories: every directed container presents a small category whose objects
are shapes and whose morphisms are positions with specified codomain; every
small category yields a directed container whose positions are outgoing
morphisms. The two constructions are mutually inverse up to a pointwise
bijection on positions, not up to definitional equality. At the morphism
level, directed-container morphisms coincide with cofunctors of the presented
categories: the two structures carry identical data, compose strictly, and the
correspondence is an isomorphism of categories. The proofs of all stated
results have been machine-checked; we record only the mathematical content.
\end{abstract}

\tableofcontents
\bigskip

%======================================================================
\section{Introduction}
%======================================================================

Directed containers equip a container $(S,P)$ with root, sub-shape, and shift
operations satisfying five laws D1--D5. Reading the data categorically---shapes
as objects, positions as morphisms, root as identity, shift as
composition---one obtains a small category. Conversely, from a small category
one builds a directed container by taking outgoing morphisms as positions.
This is the object-level dictionary of Ahman and Uustalu.

The two constructions are inverse only up to isomorphism. Starting from a
directed container $D$, the round trip replaces each position type $P(s)$ by
the total space of the fibration of positions over their codomains---each
position tagged with the shape it lands on and a proof of that landing. We
exhibit the evident pointwise bijection and leave the upgrade to an
isomorphism of directed containers for later work.

At the morphism level the situation is cleaner. A directed-container morphism
has a covariant shape map and a contravariant position map, compatible with
root, sub-shape, and shift. This is not a functor of the presented categories
(the variance on hom-sets is wrong); it is a cofunctor. Under the dictionary
the two structures are definitionally the same data, and the correspondence is
an isomorphism of categories.

\paragraph{What this note claims.}
In ordinary mathematical language: (i)~directed containers present small
categories, and small categories present directed containers, inverse up to a
canonical bijection on positions; (ii)~morphisms of directed containers
\emph{are} cofunctors of those categories, and the identification is an
isomorphism $\DCont\cong\Cof$.

%======================================================================
\section{Preliminaries}
%======================================================================

We work in a single universe $\Type$.

\begin{definition}[Directed container]\label{def:directed}
A \textbf{directed container} $D$ consists of a shape type $S$, a position
family $P$, a root $\rootpos(s)\in P(s)$, a sub-shape $s\sub p\in S$, and a
shift $p\shift q\in P(s)$ for $q\in P(s\sub p)$, subject to:
\begin{enumerate}
  \item[\textbf{D1.}] $s\sub\rootpos(s)=s$;
  \item[\textbf{D4.}] $s\sub(p\shift q)=(s\sub p)\sub q$;
  \item[\textbf{D2.}] $\rootpos(s)\shift q = (\mathrm{D1})\cdot q$
        (transport along D1);
  \item[\textbf{D3.}] $p\shift\rootpos(s\sub p)=p$;
  \item[\textbf{D5.}]
        $(p\shift q)\shift\bigl((\mathrm{D4})^{-1}\cdot q'\bigr)
        = p\shift(q\shift q')$
        (transport along D4).
\end{enumerate}
\end{definition}

\begin{definition}[Small category]\label{def:smallcat}
A \textbf{small category} $\mathcal{C}$ consists of a type of objects, a type
of morphisms $\Hom(a,b)$ between any two objects, identity morphisms, and
composition in \emph{diagrammatic} order
\[
  \mathrm{comp}:\Hom(a,b)\to\Hom(b,c)\to\Hom(a,c),
\]
subject to the usual left unit, right unit, and associativity laws.
Diagrammatic order matches the shift of a directed container: $p\shift q$
first travels along $p$, then along $q$.
\end{definition}

\begin{lemma}[Transport cancellation]\label{lem:transport-cancel}
For a family $P$ over a type $\alpha$, an equality $h:x=y$ in $\alpha$, and
$p\in P(y)$,
\begin{equation}
  h\cdot\bigl(h^{-1}\cdot p\bigr) \;=\; p.
\end{equation}
\end{lemma}

\begin{proof}
By lean.
\end{proof}

%======================================================================
\section{From directed containers to small categories}
%======================================================================

\begin{lemma}[Codomain of a composite]\label{lem:sub-shift}
Let $D$ be a directed container, and let $a,b,c\in S$ with $p\in P(a)$,
$h_p:a\sub p=b$, $q\in P(b)$, and $h_q:b\sub q=c$. Then
\begin{equation}
  a\sub\bigl(p\shift(h_p^{-1}\cdot q)\bigr) \;=\; c.
\end{equation}
\end{lemma}

\begin{proof}
By lean.
\end{proof}

\begin{remark}
This is D4 with both shape equalities made into hypotheses, so that the
subtype property of composition in the induced category is well-typed.
\end{remark}

\begin{definition}[Presented category]\label{def:to-smallcat}
Let $D$ be a directed container. The \textbf{presented small category}
$\cat(D)$ has
\begin{enumerate}
  \item objects: the shapes $S$;
  \item morphisms $a\to b$: the subtype
        $\bigl\{p\in P(a) \mid a\sub p = b\bigr\}$;
  \item identity at $a$: $\langle\rootpos(a),\,\mathrm{D1}_a\rangle$;
  \item composition of $\langle p,h_p\rangle:a\to b$ with
        $\langle q,h_q\rangle:b\to c$:
        \[
          \bigl\langle
            p\shift(h_p^{-1}\cdot q),\;
            \text{the equality of Lemma~\ref{lem:sub-shift}}
          \bigr\rangle.
        \]
\end{enumerate}
\end{definition}

\begin{theorem}[Category laws for $\cat(D)$]\label{thm:cat-laws}
The data of Definition~\ref{def:to-smallcat} satisfy the left unit, right
unit, and associativity laws of a small category. Explicitly:
\begin{enumerate}
  \item left unit uses D2 and transport cancellation
        (Lemma~\ref{lem:transport-cancel});
  \item right unit uses D3;
  \item associativity uses D5.
\end{enumerate}
In each case one substitutes the codomain equalities, after which transports
reduce definitionally (proofs of shape equalities are definitionally
irrelevant), leaving the corresponding directed-container law.
\end{theorem}

\begin{proof}
By lean.
\end{proof}

%======================================================================
\section{From small categories to directed containers}
%======================================================================

\begin{definition}[Presented directed container]\label{def:to-dcont}
Let $\mathcal{C}$ be a small category. The \textbf{presented directed
container} $D(\mathcal{C})$ has
\begin{enumerate}
  \item shapes: the objects of $\mathcal{C}$;
  \item positions over $s$: outgoing morphisms
        $\sum_{s'}\Hom_{\mathcal{C}}(s,s')$;
  \item root at $s$: $\langle s,\,\id_s\rangle$;
  \item sub-shape of $\langle s',\,f\rangle$: the codomain $s'$;
  \item shift of $\langle s',\,f\rangle$ by $\langle s'',\,g\rangle$:
        $\langle s'',\, g\circ f\rangle$ (diagrammatic composition in
        $\mathcal{C}$).
\end{enumerate}
\end{definition}

\begin{theorem}[Directed laws for $D(\mathcal{C})$]\label{thm:d-laws}
The data of Definition~\ref{def:to-dcont} satisfy D1--D5. Laws D1 and D4 hold
strictly (the codomain of a composite is the codomain of the second factor,
definitionally), so the transports in D2 and D5 vanish; the remaining laws are
exactly the category laws of $\mathcal{C}$.
\end{theorem}

\begin{proof}
By lean.
\end{proof}

%======================================================================
\section{The object-level round trip}
%======================================================================

Starting from a directed container $D$, the composite
$D\mapsto\cat(D)\mapsto D(\cat(D))$ is not definitionally $D$: its positions
over $s$ are
\[
  \sum_{s'}\bigl\{p\in P(s) \mid s\sub p = s'\bigr\},
\]
each position tagged with its codomain and a proof. The total space of this
fibration is in bijection with $P(s)$.

\begin{definition}[Round-trip position maps]\label{def:roundtrip-pos}
Let $D$ be a directed container and $s\in S$.
\begin{enumerate}
  \item \textbf{Forward.} Send $p\in P(s)$ to
        $\bigl\langle s\sub p,\; \langle p,\,\mathsf{refl}\rangle\bigr\rangle$
        in $D\bigl(\cat(D)\bigr)$'s position type at $s$.
  \item \textbf{Backward.} Send
        $\bigl\langle s',\; \langle p,\,h\rangle\bigr\rangle$ to the
        underlying position $p\in P(s)$.
\end{enumerate}
\end{definition}

\begin{theorem}[Position bijection, forward-then-back]\label{thm:pos-rt-1}
For every directed container $D$, shape $s$, and position $p\in P(s)$,
applying the forward map and then the backward map returns $p$.
\end{theorem}

\begin{proof}
By lean.
\end{proof}

\begin{theorem}[Position bijection, back-then-forward]\label{thm:pos-rt-2}
For every directed container $D$, shape $s$, and position $q$ of
$D\bigl(\cat(D)\bigr)$ at $s$, applying the backward map and then the forward
map returns $q$.
\end{theorem}

\begin{proof}
By lean.
\end{proof}

\begin{remark}[Honesty about the gap]
The pointwise bijection of Theorems~\ref{thm:pos-rt-1}
and~\ref{thm:pos-rt-2} does not by itself yield a definitional equality of
directed containers, nor automatically an isomorphism of directed containers:
one would need a notion of directed-container isomorphism that transports
root, sub-shape, and shift along the bijection. That upgrade, and the
resulting equivalence of categories $\DCont\simeq\Cat$ (Ahman--Uustalu 2016,
Theorem~3.6), remains at the paper level relative to the present formalisation.
The other direction (category to container to category) has the same flavour:
hom-sets get wrapped in a subtype over a dependent sum, equal only up to the
evident bijection.
\end{remark}

%======================================================================
\section{Morphisms as cofunctors}
%======================================================================

\begin{definition}[Directed-container morphism]\label{def:dcont-morph}
A \textbf{morphism of directed containers} $\varphi:C\Rightarrow D$ consists
of a shape map $f:\Shape(C)\to\Shape(D)$ and a contravariant position map
$f^\sharp_s:\Pos(D)\bigl(f(s)\bigr)\to\Pos(C)(s)$, satisfying:
\begin{enumerate}
  \item[\textbf{M0.}]
        $f\bigl(s\sub f^\sharp_s(h)\bigr)=f(s)\sub h$;
  \item[\textbf{M1.}]
        $f^\sharp_s\bigl(\rootpos(f(s))\bigr)=\rootpos(s)$;
  \item[\textbf{M2.}]
        $f^\sharp_s(h\shift k)
        = f^\sharp_s(h)\shift
          f^\sharp_{s\sub f^\sharp_s(h)}
            \bigl((\mathrm{M0})^{-1}\cdot k\bigr)$.
\end{enumerate}
\end{definition}

\begin{definition}[Cofunctor]\label{def:cofunctor}
A \textbf{cofunctor} $\phi:\cat(C)\to\cat(D)$ consists of an object map
$\mathrm{obj}$ and a lift
$\mathrm{lift}_s:\Pos(D)\bigl(\mathrm{obj}(s)\bigr)\to\Pos(C)(s)$,
satisfying:
\begin{enumerate}
  \item[\textbf{C0.}]
        $\mathrm{obj}\bigl(s\sub\mathrm{lift}_s(h)\bigr)
        =\mathrm{obj}(s)\sub h$;
  \item[\textbf{C1.}]
        $\mathrm{lift}_s\bigl(\rootpos(\mathrm{obj}(s))\bigr)
        =\rootpos(s)$;
  \item[\textbf{C2.}]
        $\mathrm{lift}_s(h\shift k)
        =\mathrm{lift}_s(h)\shift
          \mathrm{lift}_{s\sub\mathrm{lift}_s(h)}
            \bigl((\mathrm{C0})^{-1}\cdot k\bigr)$.
\end{enumerate}
\end{definition}

\begin{remark}
C0/C1/C2 are M0/M1/M2 read through the dictionary of $\cat(-)$: sub-shape is
codomain, root is identity, shift is composition. The two structures carry
identical data.
\end{remark}

\begin{definition}[Correspondence]\label{def:correspondence}
Every directed-container morphism $\varphi$ determines a cofunctor
$\varphi^\wedge$ by reading $(f,f^\sharp)$ as $(\mathrm{obj},\mathrm{lift})$.
Every cofunctor $\phi$ determines a directed-container morphism $\phi^\vee$ by
the inverse reading.
\end{definition}

\begin{theorem}[Hom-set bijection]\label{thm:hom-bij}
The two assignments are mutually inverse:
\begin{equation}
  \bigl(\varphi^\wedge\bigr)^\vee \;=\; \varphi,
  \qquad
  \bigl(\phi^\vee\bigr)^\wedge \;=\; \phi.
\end{equation}
Consequently there is a bijection between directed-container morphisms
$C\Rightarrow D$ and cofunctors $\cat(C)\to\cat(D)$.
\end{theorem}

\begin{proof}
By lean.
\end{proof}

%======================================================================
\section{The category of cofunctors}
%======================================================================

\begin{definition}[Identity and composition]\label{def:cof-ops}
The \textbf{identity cofunctor} on $\cat(C)$ is the identity on objects and
on positions. The \textbf{composite} of cofunctors
$\phi:\cat(C)\to\cat(D)$ and $\psi:\cat(D)\to\cat(E)$ has object map
$\mathrm{obj}_\psi\circ\mathrm{obj}_\phi$ and lift
\[
  m \;\mapsto\;
  \mathrm{lift}_\phi\bigl(s,\,
    \mathrm{lift}_\psi\bigl(\mathrm{obj}_\phi(s),\,m\bigr)
  \bigr).
\]
\end{definition}

\begin{lemma}[Composite transport]\label{lem:comp-transport}
Let $\psi:\cat(D)\to\cat(E)$, equalities $e:y=x$ in $\Shape(D)$ and
$e_2:\mathrm{obj}_\psi(x)=z$, and $k\in\Pos(E)(z)$. Then
\begin{equation}
  e^{-1}\cdot\bigl(
    \mathrm{lift}_\psi(x,\,e_2^{-1}\cdot k)
  \bigr)
  \;=\;
  \mathrm{lift}_\psi\bigl(
    y,\,
    \bigl((\mathrm{obj}_\psi\cdot e)\cdot e_2\bigr)^{-1}\cdot k
  \bigr).
\end{equation}
\end{lemma}

\begin{proof}
By lean.
\end{proof}

\begin{theorem}[Cofunctors form a strict category]\label{thm:cof-cat}
Identity and composition of cofunctors satisfy left unit, right unit, and
associativity strictly. Directed containers and cofunctors therefore form a
strict category $\Cof$.
\end{theorem}

\begin{proof}
By lean.
\end{proof}

\begin{definition}[Operations in $\DCont$]\label{def:dcont-ops}
Identity and composition of directed-container morphisms are obtained by
transporting the cofunctor operations across the hom-set bijection of
Theorem~\ref{thm:hom-bij}.
\end{definition}

\begin{theorem}[Isomorphism of categories]\label{thm:iso}
The assignment $\varphi\mapsto\varphi^\wedge$ preserves identities and
composition, is the identity on objects, and is bijective on hom-sets.
Hence it is an isomorphism of categories
\begin{equation}
  \DCont \;\cong\; \Cof.
\end{equation}
\end{theorem}

\begin{proof}
By lean.
\end{proof}

%======================================================================
\section{Discussion}
%======================================================================

\paragraph{What has been established.}
Directed containers and small categories are intertranslatable at the object
level, inverse up to a canonical position bijection. At the morphism level the
translation is tighter: directed-container morphisms and cofunctors are the
same data, and the categories they form are isomorphic.

\paragraph{Why cofunctors rather than functors.}
Container morphisms pull positions backward. Functors push morphisms forward.
The mismatch is definitional; cofunctors (object map forward, morphism action
backward) are the matching notion.

\paragraph{The remaining gap.}
Upgrading the pointwise position bijection to an equivalence of categories
$\DCont\simeq\Cat$ requires a notion of isomorphism (or equivalence) of
directed containers compatible with root, sub-shape, and shift. The present
development stops short of that packaging; the bijection itself is fully
proved.

\paragraph{Machine checking.}
Every definition and theorem statement above corresponds to a declaration in a
machine-checked formalisation. The phrase ``by lean'' indicates that the proof
has been verified by a proof assistant.

%======================================================================
\section*{Acknowledgements}
%======================================================================

The mathematical content follows Ahman and Uustalu, \emph{Directed containers
as categories} (2016), together with the cofunctor reading of
directed-container morphisms. The presentation here is a prose rendering of a
formal development intended for human review of the theorem statements.

\begin{thebibliography}{9}

\bibitem{ahman2012}
D.~Ahman, J.~Chapman, and T.~Uustalu.
\newblock When is a container a comonad?
\newblock In \emph{FoSSaCS}, volume~7213 of \emph{LNCS}, pages 74--88.
\newblock Springer, 2012.

\bibitem{ahman2016}
D.~Ahman and T.~Uustalu.
\newblock Directed containers as categories.
\newblock In \emph{MSFP}, volume~207 of \emph{EPTCS}, pages 89--98, 2016.

\bibitem{abbott2005}
M.~Abbott, T.~Altenkirch, and N.~Ghani.
\newblock Containers: constructing strictly positive types.
\newblock \emph{Theoretical Computer Science}, 342(1):3--27, 2005.

\end{thebibliography}

\end{document}
